\exercisesfor{6}

\begin{enumerate}
\def\labelenumi{\arabic{enumi}.}
\tightlist
\item
  Show that \(\exp (\mathrm{U}).\exp (\mathrm{V}).\exp (-\mathrm{U}) = \exp \left(\sum_{n = 0}^{\infty}\frac{1}{n!} (\operatorname {ad}\mathrm{U})^{n}(\mathrm{V})\right)\) . Deduce that
\end{enumerate}

\[
\mathrm{H} (\mathrm{U}, \mathrm{H} (\mathrm{V}, - \mathrm{U})) = \sum_ {n = 0} ^ {\infty} \frac {1}{n !} (\operatorname{ad} \mathrm{U}) ^ {n} (\mathrm{V}) = (\exp (\operatorname{ad} \mathrm{U})) (\mathrm{V}).
\]

\begin{enumerate}
\def\labelenumi{\arabic{enumi}.}
\setcounter{enumi}{1}
\item
  Show that \(\mathbf{H}(\mathbf{U},\mathbf{V}) = \sum_{n = 0}^{\infty}\frac{1}{n!} (\operatorname {ad}\mathbf{U})^{n}(\mathbf{H}(\mathbf{V},\mathbf{U}))\) . (Apply Exercise 1, noting that \(\mathbf{H}(\mathbf{U},\mathbf{H}(\mathbf{H}(\mathbf{V},\mathbf{U}), - \mathbf{U})) = \mathbf{H}(\mathbf{U},\mathbf{V}).\)
\item
  Let \(\mathbf{M}(\mathbf{U},\mathbf{V}) = \mathbf{H}(-\mathbf{U},\mathbf{U} + \mathbf{V})\) . Then
\end{enumerate}

\[
\exp (\mathrm{M} (\mathrm{U}, \mathrm{V})) = \exp (- \mathrm{U}). \exp (\mathrm{U} + \mathrm{V}).
\]

Let \(M_{1}(U, V)\) (resp. \(H_{1}(U, V)\) ) denote the sum of the bihomogeneous terms of the series M (resp. H) whose degree in V is 1.

\begin{enumerate}
\def\labelenumi{(\alph{enumi})}
\tightlist
\item
  Show that
\end{enumerate}

\[
\mathrm{M} _ {1} (\mathrm{U}, \mathrm{V}) = \sum_ {n = 0} ^ {\infty} \frac {(- 1) ^ {n}}{(n + 1) !} (\mathrm{adU}) ^ {n} (\mathrm{V}) = (f (\mathrm{adU})) (\mathrm{V})
\]

where \(f(\mathrm{T}) = (1 - e^{-\mathrm{T}}) / \mathrm{T}\) .

(Same argument as for the proof of Proposition 5. It can also be reduced to Proposition 5 by using the identity

\[
\exp (\mathrm{U}). \exp (\mathrm{M} (\mathrm{U}, \mathrm{V})). \exp (- \mathrm{U}) = \exp (\mathrm{H} (\mathrm{U} + \mathrm{V}, - \mathrm{U})),
\]

together with Exercise 1.)

\begin{enumerate}
\def\labelenumi{(\alph{enumi})}
\setcounter{enumi}{1}
\tightlist
\item
  Show that \(\mathrm{H}(\mathrm{U},\mathrm{M}(\mathrm{U},\mathrm{V})) = \mathrm{U} + \mathrm{V}\) . Deduce that
\end{enumerate}

\[
\mathrm{H} _ {1} (\mathrm{U}, \mathrm{M} _ {1} (\mathrm{U}, \mathrm{V})) = \mathrm{V}.
\]

\begin{enumerate}
\def\labelenumi{(\alph{enumi})}
\setcounter{enumi}{2}
\tightlist
\item
  Let \(g(\mathrm{T}) = 1 / f(\mathrm{T}) = 1 + \mathrm{T} / 2 + \sum_{n=1}^{\infty} \frac{1}{(2n)!} b_{2n}\mathrm{T}^{2n}\) , where the \(b_{2n}\) are
\end{enumerate}

the Bernoulli numbers (Functions of a Real Variable, Chapter VI, § 1, no. 4). Show that

\[
\mathrm{H} _ {1} (\mathrm{U}, \mathrm{V}) = (g (\operatorname{ad} \mathrm{U})) (\mathrm{V}) = \mathrm{V} + \frac {1}{2} [ \mathrm{U}, \mathrm{V} ] + \sum_ {n = 1} ^ {\infty} \frac {1}{(2 n) !} b _ {2 n} (\operatorname{ad} \mathrm{U}) ^ {2 n} (\mathrm{V}).
\]

Derive the first few terms of the expansion of \(H_{1}(U, V)\) :

\[
\begin{array}{r l} \mathrm{H} _ {1} (\mathrm{U}, \mathrm{V}) & = \mathrm{V} + \frac {1}{2} (\mathrm{adU}) (\mathrm{V}) + \frac {1}{12} (\mathrm{adU}) ^ {2} (\mathrm{V}) - \frac {1}{720} (\mathrm{adU}) ^ {4} (\mathrm{V}) \\ & + \frac {1}{30240} (\mathrm{adU}) ^ {6} (\mathrm{V}) - \dots \end{array}
\]

\begin{enumerate}
\def\labelenumi{(\alph{enumi})}
\setcounter{enumi}{3}
\tightlist
\item
  Show that the sum of the bihomogeneous terms of H(U, V) whose degree in U is 1 is the series:
\end{enumerate}

\[
\mathrm{U} + \frac {1}{2} [ \mathrm{U}, \mathrm{V} ] + \sum_ {n = 1} ^ {\infty} \frac {1}{(2 n) !} b _ {2 n} (\mathrm{adV}) ^ {2 n} (\mathrm{U}).
\]

(Apply (c) to H(V, U) and use Exercise 2 to pass from H(V, U) to H(U, V).) (e) Let W be another indeterminate and let X(U, V, W) be the sum of the trihomogeneous terms of H(U, V + W) whose degree in W is 1. Show that

\[
\mathrm{X} (\mathrm{U}, \mathrm{H} _ {1} (\mathrm{V}, \mathrm{W})) = \mathrm{H} _ {1} (\mathrm{H} (\mathrm{U}, \mathrm{V}), \mathrm{W}).
\]

(Use the identity \(\mathrm{H}(\mathrm{U},\mathrm{H}(\mathrm{V},\mathrm{W})) = \mathrm{H}(\mathrm{H}(\mathrm{U},\mathrm{V}),\mathrm{W}).\) ) Deduce that

\[
\mathrm{X} (\mathrm{U}, \mathrm{V}, \mathrm{W}) = (g (\operatorname{ad} \mathrm{H} (\mathrm{U}, \mathrm{V})) \circ f (\operatorname{ad} \mathrm{V})) (\mathrm{W}).
\]

¶ 4. Let X be a finite set and let \(\hat{\mathbf{L}} = \hat{\mathbf{L}}(\mathbf{X})\) . We give \(\hat{\mathbf{L}}\) the Hausdorff law of composition \((a, b) \mapsto a \uplus b\) , which we denote simply by \(a.b\) . If \(t \in \mathbf{K}\) , \(a \in \hat{\mathbf{L}}\) , we write \(a^t = t a\) .

\begin{enumerate}
\def\labelenumi{(\alph{enumi})}
\tightlist
\item
  Let \(\mathcal{H}\) be a Hall set relative to X. If \(m \in \mathcal{H}\) , let \(m_{\mathrm{L}}\) denote the corresponding element of the Lie algebra \(\mathrm{L}(\mathrm{X})\) (§2, no. 11) and \(m_{\mathcal{H}}\) the basic commutator of the group \(\hat{\mathrm{L}}\) defined by \(m\) (§5, no. 4, Remark). If \(m\) is of length \(l\) , show that \(m_{\mathcal{H}} = m_{\mathrm{L}} + \mu\) , where \(\mu\) is of filtration \(\geqslant l + 1\) (argue by induction on \(l\) ). Deduce that every element \(w\) of the group \(\hat{\mathrm{L}}\) can be written uniquely in the form
\end{enumerate}

\[
w = \prod_ {m \in \mathcal {H}} m _ {\mathcal {H}} ^ {\alpha (m)},
\]


where \(\alpha(m) \in \mathbf{K}\) and the product is convergent in the topological group \(\hat{\mathbf{L}}\) .

\begin{enumerate}
\def\labelenumi{(\alph{enumi})}
\setcounter{enumi}{1}
\item
  Let \(\mathbf{P}\) be the set of prime numbers and let \(c\) be an integer \(\geqslant 1\) . Let \(\mathbf{K} = \mathbf{Q}\) . Show that the group \(\hat{\mathbf{L}} / \hat{\mathbf{L}}_c\) is identified with the \(\mathbf{P}\) -envelope \(\mathrm{F}_{\mathbf{P}}^{c}\) of the group \(\mathrm{F}^c = \mathrm{F}(\mathrm{X}) / \mathrm{C}^c\mathrm{F}(\mathrm{X})\) , cf.~§ 5, Exercise 6.
\item
  Take X to be a set of two elements \(\{U, V\}\) . Show the existence of two families \(\alpha(m)\) , \(\beta(m)\) of rational numbers such that
\end{enumerate}

\[
\mathrm{U} + \mathrm{V} = \prod_ {m \in \mathcal {H}} m _ {\mathcal {H}} ^ {\alpha (m)}
\]

\[
[ \mathrm{U}, \mathrm{V} ] = \prod_ {m \in \mathcal {H}} m _ {\mathcal {H}} ^ {\beta (m)}.
\]

For example

\[
\begin{array}{c} \mathrm{U} + \mathrm{V} = \mathrm{U}. \mathrm{V}. (\mathrm{U}, \mathrm{V}) ^ {- \frac {1}{2}} \dots \\ {} [ \mathrm{U}, \mathrm{V} ] = (\mathrm{U}, \mathrm{V}) \dots \end{array}
\]

\begin{enumerate}
\def\labelenumi{(\alph{enumi})}
\setcounter{enumi}{3}
\tightlist
\item
  Let g be a nilpotent Lie Q-algebra with the Hausdorff group law. Let u, v be in g. Show that
\end{enumerate}

\[
u + v = \prod_ {m \in \mathcal {H}} m _ {\mathcal {H}} (u, v) ^ {\alpha (m)}
\]

\[
[ u, v ] = \prod_ {m \in \mathcal {H}} m _ {\mathcal {H}} (u, v) ^ {\beta (m)},
\]

where \(\alpha\) and \(\beta\) are the families of rational numbers defined above (``inversion of the Hausdorff formula''). (Use the continuous homomorphism \(\phi\colon\hat{L}\to\mathfrak{g}\) such that \(\phi(U)=u\) , \(\phi(V)=v\) and note that it is also a homomorphism for the Hausdorff group law.)

\begin{enumerate}
\def\labelenumi{(\alph{enumi})}
\setcounter{enumi}{4}
\tightlist
\item
  Let G be a torsion-free divisible nilpotent group (i.e.~P-torsion-free and P-divisible). If \(u \in G\) , \(t \in Q\) , \(u^{t}\) is defined ( \(\S 4\) , Exercise 16). If \(u \in G\) , \(v \in G\) , we define \(u + v\) and {[}u, v{]} by the formulae of (d) above. Show that G thus has a nilpotent Lie Q-algebra structure and that the corresponding Hausdorff law is the given group law on G (verify these assertions when G is a group \(F_{P}^{c}\) , cf.~(b)), and pass from this to the general case using homomorphisms of \(F_{P}^{r}\) into G).
\end{enumerate}

Let \(f\) be a mapping of \(G\) into a torsion-free divisible nilpotent group \(G'\) . Show that \(f\) is a (group) homomorphism if and only if it is a Lie algebra homomorphism. (The Hausdorff law therefore defines an isomorphism of the ``category'' of nilpotent Lie \(\mathbf{Q}\) -algebras onto that of torsion-free divisible nilpotent groups.)

\begin{enumerate}
\def\labelenumi{\arabic{enumi}.}
\setcounter{enumi}{4}
\tightlist
\item
  Let \(\mathfrak{g}\) be a nilpotent Lie \(\mathbf{Q}\) -algebra which we give the Hausdorff group law, denoted by \((x,y)\mapsto x.y\) .
\end{enumerate}

\begin{enumerate}
\def\labelenumi{(\alph{enumi})}
\tightlist
\item
  Show that, if \(x \in \mathfrak{g}\) , \(y \in \mathfrak{g}\) , then
\end{enumerate}

\[
x. y. x ^ {- 1} = \sum_ {n = 0} ^ {\infty} \frac {1}{n !} (\operatorname{ad} x) ^ {n} (y).
\]

(Use Exercise 2.)

\begin{enumerate}
\def\labelenumi{(\alph{enumi})}
\setcounter{enumi}{1}
\item
  Let \(\mathfrak{h}\) be a subset of \(\mathfrak{g}\) . Show that \(\mathfrak{h}\) is a Lie subalgebra (resp. an ideal) of \(\mathfrak{g}\) if and only if it is a saturated subgroup (resp. a saturated normal subgroup) of the group \(\mathfrak{g}\) . (Use the formulae of Exercise 4 (d) to pass from the group law to the Lie algebra law.)
\item
  Let \(\mathfrak{h}\) be a subgroup of the group \(\mathfrak{g}\) . Show that the \(\mathbf{P}\) -saturation of \(\mathfrak{h}\) is \(\mathbf{Q}.\mathfrak{h}\) .
\item
  Let \(\mathfrak{h}\) be a Lie subalgebra of \(\mathfrak{g}\) . Show that the centralizer (resp. normalizer) of \(\mathfrak{h}\) in the group \(\mathfrak{g}\) is the set of \(x \in \mathfrak{g}\) such that \((\operatorname{ad} x)(\mathfrak{h}) = 0\) (resp. \((\operatorname{ad} x)(\mathfrak{h}) \subset \mathfrak{h}\) ).
\item
  Show that the lower central series of the Lie algebra g coincides with that of the group g and that the associated graded Lie algebra gr(g) is the same from the ``group'' point of view as from the ``Lie algebra'' point of view.
\end{enumerate}

¶ 6. Let G be a finitely generated torsion-free nilpotent group. Show that, if n is sufficiently large, G can be embedded in the lower strict triangular group of order n over Z. (Let \((i, \overline{G})\) be a P-envelope of G, cf.~§ 5, Exercise 6; with its canonical Lie algebra structure (Exercise 4), \(\overline{G}\) is a finite-dimensional nilpotent Lie Q-algebra. Apply Ado's Theorem (Chapter 1, § 7) to \(\overline{G}\) and deduce an embedding of \(\overline{G}\) in a strict triangular group over Q. Pass from Q to Z by conjugating by a suitable integral diagonal matrix.)

